\section{Changes following external review}
\label{app:review-changes}
This appendix compares the reviewed 15 September version with the
first review revision of 17 September (local commit \artifact{f9e1956}).
The later reorganization of the paper is recorded at the end of this
appendix. Quotations and file locations refer to their respective historical
versions, before that reorganization. The first review revision's PDF is
preserved as \artifact{paper/versions/kourovka-experiment-2026-09-17-review.pdf}.
The 15 September baseline PDF is preserved as
\artifact{paper/versions/kourovka-experiment-2026-09-15.pdf}; its SHA-256 is
\begin{center}\small
\nolinkurl{ec9a5482cdfb9ee2b3937fd1980eed6448ab5f954454c9d5287d06aa8a4c3c25}.
\end{center}
The local baseline manuscript commit is \artifact{8a59409}; its public
filtered counterpart is \artifact{315e6c4ddf3b67aca8c97506f078cdc98d5413c1}.
Before/after quotations below are exact source excerpts, typeset with
this paper's macros. Their file locations are relative to \artifact{paper/};
section numbers can change when material moves. The structured change
record and source copies of both versions are in
\artifact{paper/reviews/revision-2026-09-17/}. A validation script checks
each quotation against its version. The deadline candidate ledger,
research artifacts and usage measurements have not been rewritten.

\subsection{14.72: reduced and doubled divisors}
\noindent\textbf{Review point.} First report, section 6; follow-up, section 2.\par
\noindent\textbf{Location.} \artifact{appendices/candidates/14-72.tex}.

\paragraph{Original text.}
\begin{quote}\small
It is therefore a smooth effective
Cartier divisor, of the codimension one required in
the question.
\end{quote}

\paragraph{Revised text.}
\begin{quote}\small
It is therefore a smooth effective Cartier divisor \(D\).
The equation \(x=0\) defines \(D\) set-theoretically but defines
\(2D\) as a Cartier divisor. The reduced divisor is not globally
principal on \(X\): on the zero section its class is the nonzero
order-two class of \((0,0)-\infty\) on the elliptic curve.
The restriction below supplies a global equation.
\end{quote}

\noindent\textbf{Reason.} Separate the original Cartier assertion from the stronger principal assertion.

\noindent\textbf{Validation.} Local equation x=y\textasciicircum{}2/(x\textasciicircum{}2-1); the zero-section restriction detects the nonzero two-torsion class.

\subsection{14.72: a globally principal fixed divisor}
\noindent\textbf{Review point.} Follow-up, section 2.\par
\noindent\textbf{Location.} \artifact{appendices/candidates/14-72.tex}.

\paragraph{Original text.}
\begin{quote}\small
The full fixed divisor
does not control this additional stabilizer.
\end{quote}

\paragraph{Revised text.}
\begin{quote}\small
Thus \(U\) answers the question negatively even
when ``regular hypersurface'' requires a single global equation for
the reduced fixed scheme. This restriction was obtained during revision;
it is not part of the deadline construction.
\end{quote}

\noindent\textbf{Reason.} The invariant open removes the obstruction while preserving the singular quotient. This is a post-deadline strengthening.

\noindent\textbf{Validation.} Exact ideal and invariant-localization argument in the revised proof; symbolic verification recorded in the revision audit.

\subsection{15.92: the prime is selected by the diagram}
\noindent\textbf{Review point.} Follow-up, sections 2 and 3.\par
\noindent\textbf{Location.} \artifact{appendices/candidates/15-92.tex}.

\paragraph{Original text.}
\begin{quote}\small
These cycle and parity facts are read from the construction,
not inferred from its headline generation theorem.
\end{quote}

\paragraph{Revised text.}
\begin{quote}\small
The displayed choice of \(11\) or \(13\) is the length of the unique
exceptional prime cycle in the respective diagram; it is not a free
choice. These cycle and parity facts are read from the construction,
not inferred from its headline generation theorem.
The finite transcription, diagram reconstruction and 150 cycle controls
are indexed in \artifact{results/15.92-summary.json} in the public
snapshot (Appendix~\ref{app:reproduction}); the external follow-up
reports replaying both controls.
\end{quote}

\noindent\textbf{Reason.} Explain the residue-class split and make the finite input locatable.

\noindent\textbf{Validation.} Checked the retained source-data paths and the reviewer's reported 150-row replay; the finite controls do not replace the all-size argument.

\subsection{21.106: contemporaneous work and artifact chronology}
\noindent\textbf{Review point.} First report, section 5; follow-up, sections 1 and 2.\par
\noindent\textbf{Location.} \artifact{sections/representative.tex}.

\paragraph{Original text.}
\begin{quote}\small
The experimental controls checked commutator arithmetic and finite
interpretations, while the proof above establishes the infinite
statement.
\end{quote}

\paragraph{Revised text.}
\begin{quote}\small
Contemporaneous work recorded in a public repository credits Lily Zhang
with a counterexample and proof and Evan Li with a Heisenberg
simplification; it also supplies a Lean formalisation\cite{zhang-li-concise}.
Our first local proof commit is dated 10 September 2026, 21:05:23 UTC;
the commit adding those Lean files is dated 11 September, 10:54:21 UTC.
These are artifact timestamps, not discovery or verified upload times,
and do not establish priority. We have not built that Lean development.
\end{quote}

\noindent\textbf{Reason.} Credit the parallel construction and formalisation without assigning discovery priority.

\noindent\textbf{Validation.} Inspected the public repository and commit metadata; no claim of a local Lean build.

\subsection{10.62/4.75: the prescribed relation and initial scale}
\noindent\textbf{Review point.} First author reply, item 6.\par
\noindent\textbf{Location.} \artifact{appendices/candidates/10-62.tex}.

\paragraph{Original text.}
\begin{quote}\small
All choices are made before choosing \(p\).
For a sufficiently large odd prime \(p\), use the exponent family
\(\{p\}\) and include \(w\) in the first maximal admissible family
of primitive loxodromic elements.
Proposition~6.1 then kills \(w^p\).
\end{quote}

\paragraph{Revised text.}
\begin{quote}\small
All choices are made before choosing \(p\).
Explicitly, the source's Remark~6.8 illustrates the extra condition
\(2\lambda\le L_S\delta_1\); the word \(abc\) requires
\(3\lambda\le L_S\delta_1\) instead. The permission to decrease
\(\lambda\) is the part we use. For a sufficiently large odd prime
\(p\), use the exponent family \(\{p\}\). The primitive cyclic
subgroup \(\langle abc\rangle\) is eligible for the first family;
start with it and extend to a maximal admissible selection, choosing
one representative for each eligible maximal loxodromic subgroup
up to conjugacy. Proposition~6.1 kills \(w^p\).
Neither the extra relation nor the scale \(3\lambda\) is asserted
as part of the headline statement of Theorem~6.7.
\end{quote}

\noindent\textbf{Reason.} Identify the precise use of the source construction shared by the two candidates.

\noindent\textbf{Validation.} Rechecked arXiv:2509.11958v2, Proposition 6.1 and the initialization of Theorem 6.7 with Remark 6.8.

\subsection{10.32: retain the correct lower-bound attribution}
\noindent\textbf{Review point.} First report, section 8.\par
\noindent\textbf{Location.} \artifact{appendices/candidates/10-32.tex}.

\paragraph{Original text.}
\begin{quote}\small
Rosser--Schoenfeld \cite[Theorems~4 and~9]{rosser-schoenfeld} give
\end{quote}

\paragraph{Revised text.}
\begin{quote}\small
Rosser--Schoenfeld \cite[Theorem~4, (3.14), and Theorem~9]{rosser-schoenfeld} give
\end{quote}

\noindent\textbf{Reason.} The suggested formula (3.16) is a different inequality. Add the correct formula number while retaining Theorem 4.

\noindent\textbf{Validation.} Visually checked printed page 70 of Rosser--Schoenfeld: (3.14) has denominator 2 log x and threshold 563; (3.16) has denominator log x and threshold 41.

\subsection{13.42: foundational reference and the version used}
\noindent\textbf{Review point.} First report, section 8.\par
\noindent\textbf{Location.} \artifact{appendices/candidates/13-42.tex}.

\paragraph{Original text.}
\begin{quote}\small
Here an \(A\)-group satisfies the power, conjugation, and
commuting-product axioms in Myasnikov--Remeslennikov
\cite[Definitions~1--2]{exponential-groups}.
\end{quote}

\paragraph{Revised text.}
\begin{quote}\small
The Notebook refers to Myasnikov--Remeslennikov's foundational paper
\cite{exponential-groups-1}. Here we use the power, conjugation,
commuting-product axioms and tensor-completion universal property
as restated in \cite[Definitions~1--2]{exponential-groups}.
\end{quote}

\noindent\textbf{Reason.} Add the original reference while preserving the precise definitions actually used.

\noindent\textbf{Validation.} Checked the bibliographic record of the 1994 paper and the retained definitions in Exponential groups 2.

\subsection{Usage: explain the one-cent difference}
\noindent\textbf{Review point.} First report, section 8.\par
\noindent\textbf{Location.} \artifact{sections/methods.tex}.

\paragraph{Original text.}
\begin{quote}\small
Thus the full-session estimate including compactions is USD 969.78,
of which USD 910.61 is assigned to the 48-hour research window.
\end{quote}

\paragraph{Revised text.}
\begin{quote}\small
Thus the full-session estimate including compactions is USD 969.78,
of which USD 910.61 is assigned to the 48-hour research window.
Adding the independently rounded phase totals gives USD 969.79;
the one-cent difference comes from rounding after aggregation.
\end{quote}

\noindent\textbf{Reason.} Make the effect of independent rounding explicit; no counters or rates change.

\noindent\textbf{Validation.} The exact usage outputs and their hashes are unchanged.

\subsection{Distinguish deadline review from later external review}
\noindent\textbf{Review point.} Follow-up, sections 3 and 4.\par
\noindent\textbf{Location.} \artifact{sections/discussion.tex}.

\paragraph{Original text.}
\begin{quote}\small
Finally, manuscript review is still review by the same agent that
produced the research. Repeated rereading and fresh calculations can
find errors, but they do not provide authorial independence. The
deadline archive contains zero outside reviews. Neither a preserved
PASS message nor the absence of a finding in an internal review
substitutes for specialist assessment.
\end{quote}

\paragraph{Revised text.}
\begin{quote}\small
The initial manuscript review was performed by the same agent that
produced the research. The deadline archive contains zero outside
reviews. A later external review and follow-up checked arguments and
finite inputs, identified scope and packaging issues, and withdrew
several initial objections. We preserve the correspondence in
Appendices~\ref{app:review-original}--\ref{app:reply-final}, with a
change record in Appendix~\ref{app:review-changes}. The reviewer's
scripts and logs are separately versioned; their reported outcomes
are distinguished from checks rerun during this revision.
The revised 14.72 construction makes the fixed divisor globally
principal, and the portable bundle now includes the missing small
certificate. These are later changes, not deadline achievements.
Neither a preserved PASS message nor a reviewer's acceptance
recommendation substitutes for specialist assessment or establishes
priority.
\end{quote}

\noindent\textbf{Reason.} Retain the historical zero-review count and add the subsequent review process.

\noindent\textbf{Validation.} Frozen candidate and usage ledgers remain byte-identical; the new correspondence has separate provenance.

\subsection{4.55: version the character-table indices}
\noindent\textbf{Review point.} First report, section 8.\par
\noindent\textbf{Location.} \artifact{appendices/candidates/4-55.tex}.

\paragraph{Original text.}
\begin{quote}\small
All other rows are zero in these columns.
\end{quote}

\paragraph{Revised text.}
\begin{quote}\small
These indices use \artifact{CharacterTable("3.A7")} and its 5-modular
table in GAP~4.16.1 with the character-table library \texttt{CTblLib}
version~1.3.11. All other rows are zero in these columns.
\end{quote}

\noindent\textbf{Reason.} Tie the row indices to the library used for the manuscript check.

\noindent\textbf{Validation.} The revised GAP checker prints the GAP and CTblLib versions and checks the decomposition rows, Galois orbits and projective dimensions.

\subsection{15.76(b): identify the contrasting subvariety}
\noindent\textbf{Review point.} First report, sections 4 and 8.\par
\noindent\textbf{Location.} \artifact{appendices/candidates/15-76.tex}.

\paragraph{Original text.}
\begin{quote}\small
for every proper subvariety of soluble groups: the proper metabelian
variety studied by Fernandes--Tsurkov provides a counterexample
\end{quote}

\paragraph{Revised text.}
\begin{quote}\small
for every proper subvariety of soluble groups: the proper metabelian
subvariety of exponent four studied by Fernandes--Tsurkov provides a counterexample
\end{quote}

\noindent\textbf{Reason.} State the exponent restriction on the cited counterexample.

\noindent\textbf{Validation.} Checked the retained Fernandes--Tsurkov source; the full soluble varieties remain the scope of the proposition.

\subsection{20.90: reconcile CA and CN terminology}
\noindent\textbf{Review point.} First report, section 8.\par
\noindent\textbf{Location.} \artifact{appendices/candidates/20-90.tex}.

\paragraph{Original text.}
\begin{quote}\small
It is not necessary that finite quotients be CN-groups.
\end{quote}

\paragraph{Revised text.}
\begin{quote}\small
Thus the example is a CA-group (its nonidentity centralizers are
abelian), hence a CN-group (those centralizers are pronilpotent).
It is not necessary that finite quotients be CN-groups.
\end{quote}

\noindent\textbf{Reason.} Explain why the inventory description and section title both apply.

\noindent\textbf{Validation.} The proof already establishes abelian centralizers; no additional computation is needed.

\subsection{Public filtered and local unfiltered snapshots}
\noindent\textbf{Review point.} Follow-up, section 2; repository clarification from the organizers.\par
\noindent\textbf{Location.} \artifact{appendices/reproduction.tex}.

\paragraph{Original text.}
\begin{quote}\small
The configured repository address is
\url{https://github.com/JUrban/kour1}.
Anonymous access could not be established during manuscript
preparation: unauthenticated API requests returned HTTP~404.
A local Git query obtained a remote head different from
the frozen experimental revision. These observations
do not establish that either revision is anonymously
available or that the two histories are equivalent.
Readers need the frozen archive identified in
Appendix~\ref{app:protocol} to reproduce repository-based
measurements. The mathematical exposition in the paper
does not require executing that archive.


\end{quote}

\paragraph{Revised text.}
\begin{quote}\small
The public repository is \url{https://github.com/JUrban/kour1}.
Its research snapshot is
\begin{center}\small
\href{https://github.com/JUrban/kour1/tree/bbfac1ac810117da37f01716d4594dc2b2e96980}{\texttt{bbfac1ac810117da37f01716d4594dc2b2e96980}}.
\end{center}
The owner filtered the history with
\artifact{--strip-blobs-bigger-than 90M}. Its commit identifiers therefore
differ from those of the original unfiltered archive, including
\snapshot{} in Appendix~\ref{app:protocol}. The unfiltered identifiers
in the historical timelines are provenance identifiers, not public
checkout instructions. On 17 September we compared the two final research
trees: all 4,146 retained file blobs have identical Git object IDs;
the sole omitted blob at that snapshot is the 1.27\,GB compressed
\kn{20.100} certificate discussed below. This comparison does not
assert equivalence of the full histories.
\artifact{paper/data/public-artifacts.json} records both snapshots,
the omitted file and byte comparisons of selected public artifacts.
The public baseline manuscript commit is
\artifact{315e6c4ddf3b67aca8c97506f078cdc98d5413c1}; its PDF is byte-identical
to the reviewed version preserved with this revision.

Recomputing Git measurements by the original scripts requires the
unfiltered history, or an explicit mapping to the rewritten commits.
The included frozen measurements retain their original identifiers.
Typesetting and the portable carpet checks require neither history.


\end{quote}

\noindent\textbf{Reason.} Replace the obsolete access failure and separate public checkout references from historical local hashes.

\noindent\textbf{Validation.} Compared all 4,146 retained Git blob IDs; byte-checked five selected public files and the baseline PDF.

\subsection{A dated retrieval index}
\noindent\textbf{Review point.} First author reply, item 4; follow-up, section 1.\par
\noindent\textbf{Location.} \artifact{appendices/reproduction.tex}.

\paragraph{Original text.}
\begin{quote}\small
The raw session inputs remain separate from both source bundles.
Their hashes and sizes identify the evidence for an authorized holder;
the included projections are not substitutes for the full trace.
\end{quote}

\paragraph{Revised text.}
\begin{quote}\small
The raw session inputs remain separate from both source bundles.
Their hashes and sizes identify the evidence for an authorized holder;
the included projections are not substitutes for the full trace.
The additional \artifact{extract_retrievals.py} parser exports
\artifact{retrieval-observations.jsonl}, \artifact{retrieval-urls.csv}
and \artifact{retrieval-metrics.json}. These index 1,324 completed
web-tool observations (1,293 during research and 31 during manuscript
preparation), preserving timestamps, action types, queries and returned
links. They distinguish a requested page action from a URL appearing
among search results. Shell downloads are outside this index; neither
kind of record establishes that a source was used in an argument.
The exact URL of the contemporaneous 21.106 repository has zero literal
occurrences in the supplied raw JSONL. That narrow check cannot exclude
exposure through other URLs, sources or training data, and the derived
index alone cannot independently establish an absence in the raw trace.
\end{quote}

\noindent\textbf{Reason.} Expose structured source-access observations without treating retrieval as mathematical use.

\noindent\textbf{Validation.} Reparsed the identified 540,278,152-byte input; all output records come from completed web observations. Raw reasoning and directive bodies are excluded.

\subsection{20.100: name the undistributed input}
\noindent\textbf{Review point.} Follow-up, section 2; repository clarification from the organizers.\par
\noindent\textbf{Location.} \artifact{appendices/reproduction.tex}.

\paragraph{Original text.}
\begin{quote}\small
Large computations need substantial time and storage.
In particular, the \(n=7\) certificate for \kn{20.100}
is distinct from smaller diagnostic and intermediate
files. The retained report and manifest specify the
complete input and all completed checker shards.
The manuscript's review does not silently replace that
historical verification with a fresh abbreviated run.
\end{quote}

\paragraph{Revised text.}
\begin{quote}\small
The \(n=7\) certificate for \kn{20.100},
\artifact{results/20.100-n7-certificate.g.gz}, is not distributed in the
filtered public repository or either paper source bundle. Its compressed
size is 1,271,256,410 bytes and its SHA-256 digest is
\begin{center}\small
\nolinkurl{e1b116fa80a4b77a03256d69480b996318522c227ef48e6be7fdf62cbd8bcd8c}.
\end{center}
The public archive retains the checker, generator and shard-verifier
logs, and final report. They document the historical verification;
without the certificate they do not permit a fresh complete replay.
The early \artifact{20.100-n7-integrity.json} receipt predates completion
of verification and must not be read as the final status.
No separate large-file deposit or new \(n=7\) verification is claimed.
A later deposit could supply this missing input without rewriting the
research record.
\end{quote}

\noindent\textbf{Reason.} State exactly what the public archive does and does not permit a reader to replay.

\noindent\textbf{Validation.} Public tree comparison confirms the omitted 1,271,256,410-byte blob; the retained generator and final verification records identify it.

\subsection{Inventory: scope, dependencies and prior-work status}
\noindent\textbf{Review point.} Follow-up, sections 1 and 5.\par
\noindent\textbf{Location.} \artifact{appendices/inventory.tex}.

\paragraph{Original text.}
\begin{quote}\small
Several entries share imported ingredients. The mathematical discussions
give the relevant scope and dependencies; the frozen ledger supplies the
original statement and evidence paths.
\end{quote}

\paragraph{Revised text.}
\begin{quote}\small
The additional review column records selected scope qualifications and
imported dependencies; an em dash means no additional note here, not
that the proof has no imports. Priority remains unestablished for all
46 candidates.
\end{quote}

\noindent\textbf{Reason.} Preserve all historical descriptions and covered parts while adding neutral qualifications. A separate table makes 21.132, 21.42 and 20.33 visible under prior-work status.

\noindent\textbf{Validation.} The generator checks exactly the same 46 keys; annotations are stored separately. No count is inferred from proof length or presumed proposer intent.

\subsection{19.56 as a fourth main example}
\noindent\textbf{Review point.} First author reply, item 6.\par
\noindent\textbf{Location.} \artifact{sections/representative.tex}.

\paragraph{Original text.}
\begin{quote}\small
The following arguments illustrate three different kinds of output.
\end{quote}

\paragraph{Revised text.}
\begin{quote}\small
The following arguments illustrate four different kinds of output.
\end{quote}

\noindent\textbf{Reason.} Move the existing 19.56 exposition from the candidate appendix into Section 4. The hyperfocal, minimal-simple and Baer--Suzuki dependencies remain explicit.

\noindent\textbf{Validation.} The 19.56 source is byte-identical; generated input lists include it once, giving four main examples and 42 further candidate expositions.

\subsection{Portable carpet bundle: supply the fourth checker input}
\noindent\textbf{Review point.} Follow-up, sections 2 and 5.\par
\noindent\textbf{Location.} \artifact{appendices/reproduction.tex}.

\paragraph{Original text.}
\begin{quote}\small
Their README gives three standard-Python
commands, which require neither GAP nor network access. A manifest
binds the 23 original files to their byte-identical research versions.
\end{quote}

\paragraph{Revised text.}
\begin{quote}\small
Their README gives four standard-Python
commands, which require neither GAP nor network access. A manifest
binds the 24 original files to their byte-identical research versions.
\end{quote}

\noindent\textbf{Reason.} Add results/19.62-g2-monomial-certificates.jsonl (2,194,486 bytes) to the ancillary directory and both source bundles; retain and document the fourth checker. The first reply overstated the completeness of the bundled data.

\noindent\textbf{Validation.} Check all four commands after extraction outside the repository, and verify each original file against the ancillary manifest.


\subsection{Withdrawn objections and retained scope}
The follow-up withdraws the algebraic-closedness objection to 16.28(a):
the original field was already the algebraic closure of \(\F_5(t)\).
That proof is unchanged. We retain characteristic five in part (b),
without asserting the identical conjugacy-class argument over every odd
characteristic. The suggested unrestricted extension is unnecessary;
in characteristic three its trace-one class does not have the same
semisimple description.

For 12.40, the follow-up accepts that nonexistence of any finite-valued
bound is a negative answer to the printed question. No proof change is
needed. The estimated count of ``roughly seven'' unintended readings,
the speculation about proposer intent, and the priority inference from
artifact dates are withdrawn in the follow-up and are not adopted here.
Proof length and speed of formalisation are not measures of discovery
difficulty. The original reports retain those earlier statements solely
as part of the correspondence record.

For 15.65, the source-index discrepancy discussed in its proof remains
explicit. The revision does not claim a new resolution of that discrepancy.
For 21.114, the external follow-up checks the retained order-512 inputs,
their central quotients and derived series. The order-8,192 construction
in the paper still uses the subsequent central-product argument; the
report is not described as a direct enumeration of that final group.
The 46 historical candidate classifications and their unresolved priority
status remain unchanged.

\subsection{Correspondence, packaging and validation records}
The four documents reproduced next were added to the manuscript without
rewriting their words. The lightly updated first report is retained as
\artifact{external-reviews/review-update1.md}, with its exact diff in
\artifact{external-reviews/review-update1.diff}. The final reply corrects
our first reply's claims about bundle completeness and large-certificate
availability. The digest manifest binds the raw documents and their
rendered counterparts. Markdown headings and tables have print formatting;
this does not change the review text or its historical references.

Both source bundles include the correspondence and the missing small
certificate. The full source additionally includes the retrieval index,
structured change record, baseline excerpts and PDF, and revision checks.
The build receipt, structure audit and portable-validation receipt bind
the resulting sources, PDF and ancillary inputs. The latter records all
four documented ancillary executions after archive extraction, and a
standalone typesetting comparison. The supplemental symbolic check of
the 14.72 restriction is an exact algebra check, not a formalisation of
the geometric proof. None of these checks is a fresh replay of the
undistributed 20.100 certificate or of all the reviewer's programs.

\subsection{Subsequent editorial reorganization and reviewer credit}
After the first review revision, the organizers identified the external
reviewer as Anthropic's Claude Opus[1m], model
\artifact{claude-opus-5[1m]}, with \texttt{xhigh} reasoning through Claude
Code. This identification is now credited on the opening page and in
Section~\ref{sec:external-review}; it is not inserted into the historical
reports or replies.

The abstract and introduction now summarize both the mathematical results
and the independent review. A new early section describes Claude's checks,
the exchange with Codex, and its confirmations and corrections. The
original instruction is summarized in the account of the research process.
Detailed configuration, accounting, chronology and study-design material
are in Appendices~\ref{app:protocol}--\ref{app:evaluation}.

The former mathematical appendices now appear in
Section~\ref{app:inventory} and
Sections~\ref{app:candidates}--\ref{app:prior}. The 42 further candidate
proofs are grouped by subject, and the opening table points to selected
results. Their hypotheses, arguments, imported dependencies and historical
classifications are preserved. The contents lists major sections so that
technical indexing does not delay the mathematical narrative.

The 18 before/after pairs above remain tied to the actual first review
revision, rather than being silently replaced by later wording.
Source copies bind both historical sides; the preserved PDF fixes the
then-current layout. The subsequent changes to organization and exposition
are recorded here and in the Git revision, with a second reading pass and
validation notes in \artifact{paper/reviews/polish-2026-09-17/}.

In a final editorial pass, the full inventory moved directly after the
introduction, before the proof sections, with updated reading routes.
A title-page footnote records Michael Kinyon's research visit to the
AI4REASON institute. Urban's title-page footnote credits ERC NextReason
(Grant Agreement No.\ 101200949) and links to the institute's sponsors;
it replaces the separate acknowledgements paragraph.

The abstract was subsequently shortened by removing a repeated assessment
of review evidence and the deadline count. The introduction now states
the result categories and the nine partial-result questions explicitly,
and places the deadline and requested local resource limits alongside
the experiment's objective. Its motivation relates this evaluation on
open problems to the MPTPChallenge and timed human mathematical contests.
